\section{Formal Verification Methodology}
\label{sec:formal-verification}

This section uses a slightly changed search for the maximal sublevel set $\mathcal{V}_\rho$ as 
proposed in \cite{yangLyapunovstableNeuralControl2024}. For saving compute during 
non-certificable lyapunov candidates, the algorithm includes limits for scaling and bisection. 
This leads to the following algorithm:

\begin{algorithm}[htpb]
    \caption{Certified Value Search via Scaling and Bisection}
    \label{alg:search-and-bisect}
    \SetAlgoLined
    \KwIn{$\rho_0, \rho_{\min}$, evaluator $f$, scale $s > 1$, tol $\varepsilon$, limits $N_{\text{scale}}, N_{\text{bisect}}$}
    \KwOut{Max certified value $\rho^*$}

    \BlankLine
    $\rho \leftarrow \max(\rho_0, \rho_{\min})$\;

    \BlankLine
    \tcp{Phase 1: Scaling}
    \eIf{$f(\rho)$}{
        \For{$k \leftarrow 1$ \KwTo $N_{\text{scale}}$}{
            \If{$\neg f(\rho)$}{\textbf{break}\;}
            $\rho_{\text{lo}} \leftarrow \rho;\quad \rho \leftarrow \rho \cdot s$\;
        }
        $\rho_{\text{up}} \leftarrow \rho$\;
    }{
        \For{$k \leftarrow 1$ \KwTo $N_{\text{scale}}$}{
            \If{$f(\rho) \lor \rho = \rho_{\min}$}{\textbf{break}\;}
            $\rho_{\text{up}} \leftarrow \rho;\quad \rho \leftarrow \max(\rho_{\min}, \rho / s)$\;
        }
        \If{$\neg f(\rho)$}{\KwRet $0$\ \tcp*{No valid value $\ge \rho_{\min}$}}
        $\rho_{\text{lo}} \leftarrow \rho$\;
    }

    \BlankLine
    \tcp{Phase 2: Bisection}
    \For{$k \leftarrow 1$ \KwTo $N_{\text{bisect}}$}{
        \If{$\rho_{\text{up}} - \rho_{\text{lo}} \le \varepsilon$}{\textbf{break}\;}
        $\rho_{\text{mid}} \leftarrow \frac{1}{2}(\rho_{\text{lo}} + \rho_{\text{up}})$\;
        \eIf{$f(\rho_{\text{mid}})$}{
            $\rho_{\text{lo}} \leftarrow \rho_{\text{mid}}$\;
        }{
            $\rho_{\text{up}} \leftarrow \rho_{\text{mid}}$\;
        }
    }

    \BlankLine
    \KwRet $\rho_{\text{lo}}$\;
\end{algorithm}

Regarding evaluation, the $\alpha\beta$-Crown framework is applied on a certification formulation, 
described as

\begin{equation}
    \label{eq:outside-certification}
    \Phi_{\text{sub}}(x; \rho) \equiv 
        V_\theta(x) > \rho
\end{equation}

\begin{equation}
    \label{eq:positive-certification}
    \Phi_{\text{pos}}(x) \equiv 
        V_\theta(x) > 0
\end{equation}

\begin{equation}
    \label{eq:decrease-certification}
    \Phi_{\text{dec}}(x) \equiv 
        (1 - \kappa) V_\theta(x) - V_\theta(f_d(x, \pi_\theta(x))) > 0
\end{equation}

\begin{equation}
    \label{eq:invariance-certification}
    \Phi_{\text{inv}}(x) \equiv \left(
        f(x, \pi_\theta(x)) > \underbar{x}_C
    \right) \land \left(
        f(x, \pi_\theta(x)) < \bar{x}_C
    \right)
\end{equation}

\begin{equation}
    \label{eq:condition-certification}
    \Phi_{\text{cond}}(x) \equiv 
        \Phi_{\text{pos}}(x) 
        \land \Phi_{\text{dec}}(x) 
        \land \Phi_{\text{inv}}(x)
\end{equation}


\begin{equation}
    \label{eq:total-certification-formulation}
    \Phi(x; \rho) \equiv \Phi_{\text{sub}}(x; \rho) \lor \Phi_{\text{cond}}(x)
\end{equation}
